\bta{2004}


\begin{enumerate}
	%\renewcommand{\labelenumi}{\arabic{enumi}.}
	% A(\Alph) a(\alph) I(\Roman) i(\roman) 1(\arabic)
	%设定全局标号series=example	%引用全局变量resume=example
	%[topsep=-0.3em,parsep=-0.3em,itemsep=-0.3em,partopsep=-0.3em]
	%可使用leftmargin调整列表环境左边的空白长度 [leftmargin=0em]
	\item
	填空题 (本题共 5 题，每小题 5 分，满分 25 分)
	\begin{enumerate}
		%\renewcommand{\labelenumi}{\arabic{enumi}.}
		% A(\Alph) a(\alph) I(\Roman) i(\roman) 1(\arabic)
		%设定全局标号series=example	%引用全局变量resume=example
		%[topsep=-0.3em,parsep=-0.3em,itemsep=-0.3em,partopsep=-0.3em]
		%可使用leftmargin调整列表环境左边的空白长度 [leftmargin=0em]
		\item
	$\lim\limits _{n \rightarrow \infty} \sqrt[n]{\sin 1+\sin \frac{1}{2}+\cdots+\sin \frac{1}{n}}=$ \tk{$ 1 $} 。
		
		
		\item
		设 $y-\epsilon \sin y=x$ (常数 $\epsilon \in(0,1))$, 则 $\frac{d^{2} y}{d x^{2}}=$ \tk{$ -\frac{\epsilon \sin y}{(1-\epsilon \cos y)^{3}} $} 。
		
		
		\item
		积分 $\int_{0}^{+\infty} \frac{\ln \left(1+x^{2}\right)}{x^{\alpha}} d x$ 的收敛域为 \tk{$ 1< \alpha <3 $} 。
		
		
		\item
		 曲面 $z=\arctan \frac{y}{x}$ 在点 $\left(1,1, \frac{\pi}{4}\right)$ 处的切平面方程为 \tk{$ x-y+2z=\frac{\pi}{2} $} 。
		
		
		\item
		微分方程 $y^{\prime \prime}-3 y^{\prime}+2 y=\cos x$ 的通解为 \tk{$ y= \frac{1}{10}(\cos x -3\sin x )+C_{1}e^{x}+C_{2}e^{2x} $} 。
		
		
		
		
	\end{enumerate}

\banswer{
\begin{enumerate}
	%\renewcommand{\labelenumi}{\arabic{enumi}.}
	\item
	$1$。记 $S_{n}=\sum\limits_{k=1}^{n}\sin\frac1k$。当 $k\geqslant1$ 时 $0<\frac1k\leqslant1<\frac{\pi}{2}$，
	故 $0<\sin\frac1k\leqslant\sin1<1$，从而
	\[
	\sin1\leqslant S_{n}\leqslant n\sin1\leqslant n ,
	\]
	于是 $(\sin1)^{\frac1n}\leqslant\sqrt[n]{S_{n}}\leqslant n^{\frac1n}$。
	因 $(\sin1)^{1/n}\to1$、$n^{1/n}\to1$，由夹逼准则知原极限为 $1$。

	\item
	$-\frac{\epsilon\sin y}{(1-\epsilon\cos y)^{3}}$。对 $y-\epsilon\sin y=x$ 两端关于 $x$ 求导得
	$(1-\epsilon\cos y)\frac{dy}{dx}=1$；因 $\epsilon\in(0,1)$ 有 $1-\epsilon\cos y>0$，故
	$\frac{dy}{dx}=\frac{1}{1-\epsilon\cos y}$。再对 $x$ 求导：
	\begin{align*}
	\frac{d^{2}y}{dx^{2}}&=-\frac{1}{(1-\epsilon\cos y)^{2}}\cdot\epsilon\sin y\cdot\frac{dy}{dx}\\
	&=-\frac{\epsilon\sin y}{(1-\epsilon\cos y)^{3}} .
	\end{align*}

	\item
	$1<\alpha<3$。被积函数非负，两个端点都要考察。
	$x\to0^{+}$ 时 $\ln\left(1+x^{2}\right)\sim x^{2}$，被积函数 $\sim x^{2-\alpha}$，
	故 $\int_{0}^{1}$ 收敛当且仅当 $2-\alpha>-1$，即 $\alpha<3$；
	$x\to+\infty$ 时 $\ln\left(1+x^{2}\right)\sim2\ln x$，被积函数 $\sim\frac{2\ln x}{x^{\alpha}}$，
	故 $\int_{1}^{+\infty}$ 收敛当且仅当 $\alpha>1$。两者同时成立即得收敛域 $1<\alpha<3$。

	\item
	$x-y+2z=\frac{\pi}{2}$。由 $z=\arctan\frac{y}{x}$ 得
	\[
	z_{x}=\frac{-y}{x^{2}+y^{2}},\qquad z_{y}=\frac{x}{x^{2}+y^{2}} ,
	\]
	在 $(1,1)$ 处 $z_{x}=-\frac12$，$z_{y}=\frac12$。故切平面为
	$z-\frac{\pi}{4}=-\frac12(x-1)+\frac12(y-1)$，整理得 $x-y+2z=\frac{\pi}{2}$。

	\item
	$y=\frac{1}{10}(\cos x-3\sin x)+C_{1}e^{x}+C_{2}e^{2x}$。
	特征方程 $r^{2}-3r+2=0$ 的根为 $r_{1}=1$、$r_{2}=2$，齐次通解为 $C_{1}e^{x}+C_{2}e^{2x}$。
	因 $\pm i$ 不是特征根，设特解 $y^{*}=A\cos x+B\sin x$，代入原方程得
	\[
	(A-3B)\cos x+(3A+B)\sin x=\cos x ,
	\]
	故 $A-3B=1$，$3A+B=0$，解得 $A=\frac{1}{10}$，$B=-\frac{3}{10}$，
	即 $y^{*}=\frac{1}{10}(\cos x-3\sin x)$。
\end{enumerate}
}

\item
单项选择题 (本题共 5 题，每小题 5，满分 25 分)
\begin{enumerate}
	%\renewcommand{\labelenumi}{\arabic{enumi}.}
	% A(\Alph) a(\alph) I(\Roman) i(\roman) 1(\arabic)
	%设定全局标号series=example	%引用全局变量resume=example
	%[topsep=-0.3em,parsep=-0.3em,itemsep=-0.3em,partopsep=-0.3em]
	%可使用leftmargin调整列表环境左边的空白长度 [leftmargin=0em]
	\item
设 $S$ 为球面 $x^{2}+y^{2}+z^{2}=R^{2}$ 外侧，则 $\oiint_{S} x^{2} d y d z+y^{2} d z d x+z^{2} d x d y=$ \xzanswer{A}

	
\fourchoices
{$ 0 $}
{$\pi R^{4}$}
{$2 \pi R^{4}$}
{$4 \pi R^{4}$}

	
	
	\item
	曲线 $y=\sqrt{x^{2}+1}-x-1$ 的渐进线条数为 \xzanswer{C} 
	
	
\fourchoices
{$ 0 $}
{$ 1 $}
{$ 2 $}
{$ 3 $}


\item 
给定严格递增数列 $\left\{A_{n}\right\}$, 且 $A_{1}=a, \lim\limits _{n \rightarrow \infty} A_{n}=+\infty$, 函数 $f(x)$ 在 $[a,+\infty)$ 上连
续且非负，则积分 $\int_{a}^{+\infty} f(x) d x$ 收敛是级数 $\sum\limits_{n=1}^{\infty} \int_{A_{n}}^{A_{n+1}} f(x) d x$ 收敛的 \xzanswer{C}



\fourchoices
{充分条件但非必要条件}
{必要条件但非充分条件}
{充分必要条件}
{既非充分条件也非必要条件}

\item 
如果级数 $\sum\limits_{n=2}^{\infty} \ln \left[1+\frac{(-1)^{n}}{n^{p}}\right](p>0)$ 条件收敛，则 \xzanswer{D} 


\fourchoices
{$0<p \leqslant 1$}
{$p>1$}
{$\frac{1}{3}<p \leqslant 1$}
{$\frac{1}{2}<p \leqslant 1$}


\item 
设$f(x, y)=\left\{\begin{array}{ll}\left(x^{2}+y^{2}\right) \sin \frac{1}{x^{2}+y^{2}} & (x, y) \neq(0,0) \\ 0 & (x, y)=(0,0)\end{array}\right.$，则下列选项正确的是 \xzanswer{D} 

\fourchoices
{$f(x, y)$ 在 $(0,0)$ 处不可微， $\frac{\partial f}{\partial x}, \frac{\partial f}{\partial y}$ 在 $(0,0)$ 处连续}
{$f(x, y)$ 在 $(0,0)$ 处不可微， $\frac{\partial f}{\partial x}, \frac{\partial f}{\partial y}$ 在 $(0,0)$ 处不连续}
{$f(x, y)$ 在 $(0,0)$ 处可微， $\frac{\partial f}{\partial x}, \frac{\partial f}{\partial y}$ 在 $(0,0)$ 处连续}
{$f(x, y)$ 在 $(0,0)$ 处可微， $\frac{\partial f}{\partial x}, \frac{\partial f}{\partial y}$ 在 $(0,0)$ 处不连续}

	
	
\end{enumerate}


\banswer{
\begin{enumerate}
	%\renewcommand{\labelenumi}{\arabic{enumi}.}
	\item
	\textbf{A}。记 $\Omega$ 为球体 $x^{2}+y^{2}+z^{2}\leqslant R^{2}$。$S$ 取外侧即 $\Omega$ 的正向边界，
	由高斯公式
	\begin{align*}
	\oiint_{S}x^{2}dydz+y^{2}dzdx+z^{2}dxdy
	&=\iiint_{\Omega}(2x+2y+2z)\,dV\\
	&=0 ,
	\end{align*}
	因为 $\Omega$ 关于三个坐标平面都对称，而 $x$、$y$、$z$ 分别关于相应变量为奇函数。

	\par\vspace{1mm}\noindent
	\textbf{注：}{\kaishu
	原题面第三项印作 $z\,dxdy$。若按此读，由高斯公式得
	$\iiint_{\Omega}(2x+2y+1)dV=\frac43\pi R^{3}$，并不在四个选项之中；
	而四个选项的量阶均为 $R^{4}$，与被积三项全为二次时一致，
	可见应为 $z^{2}dxdy$，已据此更正题面。
	}\par

	\item
	\textbf{C}。函数在 $\mathbb{R}$ 上连续，无铅直渐近线。
	$x\to+\infty$ 时 $\sqrt{x^{2}+1}-x=\frac{1}{\sqrt{x^{2}+1}+x}\to0$，故 $y\to-1$，
	得水平渐近线 $y=-1$。
	$x\to-\infty$ 时 $\sqrt{x^{2}+1}=-x\sqrt{1+\frac{1}{x^{2}}}=-x-\frac{1}{2x}+o\!\left(\frac1x\right)$，
	故
	\[
	y=-2x-1-\frac{1}{2x}+o\!\left(\frac1x\right) ,
	\]
	得斜渐近线 $y=-2x-1$。共 $2$ 条。

	\item
	\textbf{C}。因 $f\geqslant0$，令 $F(X)=\int_{a}^{X}f(x)dx$，则 $F$ 单调不减。级数的部分和
	\[
	\sum_{n=1}^{N}\int_{A_{n}}^{A_{n+1}}f(x)dx=\int_{A_{1}}^{A_{N+1}}f(x)dx=F\left(A_{N+1}\right) .
	\]
	若积分收敛，则 $F$ 有极限，部分和收敛，故为充分条件；
	反之若部分和收敛，因 $A_{N}\uparrow+\infty$ 且 $F$ 单调不减有界，
	故 $\lim\limits_{X\to+\infty}F(X)$ 存在，积分收敛，故也是必要条件。二者等价。

	\par\vspace{1mm}\noindent
	\textbf{注：}{\kaishu
	原题面级数印作 $\sum\limits_{n=1}^{\infty}\int_{A_{n}}^{A_{n+1}}dx$，该式恒等于
	$\lim\limits_{N\to\infty}A_{N+1}-a=+\infty$ 必定发散，题即无意义；
	显系漏印被积函数 $f(x)$，已补上。
	}\par

	\item
	\textbf{D}。记 $u_{n}=\ln\left[1+\frac{(-1)^{n}}{n^{p}}\right]$。由
	$\ln(1+t)=t-\frac{t^{2}}{2}+O\left(t^{3}\right)$ 得
	\[
	u_{n}=\frac{(-1)^{n}}{n^{p}}-\frac{1}{2n^{2p}}+O\left(n^{-3p}\right) .
	\]
	其中 $\sum\frac{(-1)^{n}}{n^{p}}$ 对一切 $p>0$ 收敛（莱布尼茨），
	而 $\sum\frac{1}{2n^{2p}}$ 收敛当且仅当 $2p>1$。故级数收敛当且仅当 $p>\frac12$。
	又 $\left|u_{n}\right|\sim\frac{1}{n^{p}}$，故绝对收敛当且仅当 $p>1$。
	所以条件收敛当且仅当 $\frac12<p\leqslant1$。

	\par\vspace{1mm}\noindent
	\textbf{注：}{\kaishu
	原题面选项 A 印作 “$0<\leqslant1$”、选项 D 印作 “$\frac12<p<\leqslant1$”，
	均漏字或衍字，已分别更正为 $0<p\leqslant1$ 与 $\frac12<p\leqslant1$。
	}\par

	\item
	\textbf{D}。沿坐标轴 $f(x,0)=x^{2}\sin\frac{1}{x^{2}}$，故
	\[
	f_{x}(0,0)=\lim_{x\to0}\frac{x^{2}\sin\frac{1}{x^{2}}}{x}=\lim_{x\to0}x\sin\frac{1}{x^{2}}=0 ,
	\]
	同理 $f_{y}(0,0)=0$。记 $\rho=\sqrt{x^{2}+y^{2}}$，则
	\[
	\frac{f(x,y)-f(0,0)-0\cdot x-0\cdot y}{\rho}=\rho\sin\frac{1}{\rho^{2}}\longrightarrow0 ,
	\]
	故 $f$ 在 $(0,0)$ 处可微。但当 $(x,y)\neq(0,0)$ 时
	\[
	f_{x}=2x\sin\frac{1}{\rho^{2}}-\frac{2x}{\rho^{2}}\cos\frac{1}{\rho^{2}} ,
	\]
	沿 $y=0$ 有 $f_{x}(x,0)=2x\sin\frac{1}{x^{2}}-\frac{2}{x}\cos\frac{1}{x^{2}}$，
	当 $x\to0$ 时无界振荡（例如取 $x_{k}=\frac{1}{\sqrt{2k\pi}}$ 得 $-2\sqrt{2k\pi}\to-\infty$），
	故 $f_{x}$ 在 $(0,0)$ 处不连续，$f_{y}$ 同理。即可微而偏导数不连续。
\end{enumerate}
}


\item
(本题共 5 题，每小题 8 分，满分 40 分)
\begin{enumerate}
	%\renewcommand{\labelenumi}{\arabic{enumi}.}
	% A(\Alph) a(\alph) I(\Roman) i(\roman) 1(\arabic)
	%设定全局标号series=example	%引用全局变量resume=example
	%[topsep=-0.3em,parsep=-0.3em,itemsep=-0.3em,partopsep=-0.3em]
	%可使用leftmargin调整列表环境左边的空白长度 [leftmargin=0em]
	\item
	计算极限 $\lim\limits _{x \rightarrow 0} \frac{\int_{0}^{\tan x} t(\tan t-t) d t}{\int_{0}^{\sin ^{2} x} \sin ^{\frac{3}{2}} t d t}$。
	\item 
计算积分 $\int_{0}^{\ln 2} \sqrt{e^{x}-1} d x$。
	
	
	\item
	利用欧拉积分计算 $\int_{0}^{\frac{\pi}{2}}(\tan x)^{\frac{2}{3}} d x$。
	
	
	



\item
利用 Stokes 公式计算
\[ 
\oint_{L}(y-z) d x+(z-x) d y+(x-y) d z
 \]
 其中$L:\left\{\begin{array}{l}x^{2}+y^{2}+z^{2}=a^{2} \\ x+y+z=0\end{array}\right.$ $ (a>0) $，从 $x$ 轴正向看 $L$ 为逆时针走向。


\item
 设 $a, b>0$ 。 证明：当 $y>x>0$ 时，有 $\left(a^{x}+b^{x}\right)^{\frac{1}{x}}>\left(a^{y}+b^{y}\right)^{\frac{1}{y}}$。

\end{enumerate}


\banswer{
	\begin{enumerate}
		%\renewcommand{\labelenumi}{\arabic{enumi}.}
		% A(\Alph) a(\alph) I(\Roman) i(\roman) 1(\arabic)
		%设定全局标号series=example	%引用全局变量resume=example
		%[topsep=-0.3em,parsep=-0.3em,itemsep=-0.3em,partopsep=-0.3em]
		%可使用leftmargin调整列表环境左边的空白长度 [leftmargin=0em]
		\item
		$  \frac{ 1 }{ 6 }  $
		\item 
		$ 2-\frac{\pi}{2} $		
		\item 
		$ \pi $
		\item 
		$ -2\sqrt{3}\pi a^{2} $
		\item 
		\textbf{证} 令 $p=\dfrac{y}{x}>1$，记 $A=a^{x}>0$，$B=b^{x}>0$，则 $a^{y}=A^{p}$，$b^{y}=B^{p}$。只需证 $(A+B)^{p}>A^{p}+B^{p}$。两边除以 $B^{p}$，令 $t=\dfrac{A}{B}>0$，即证 $(1+t)^{p}>1+t^{p}$。令
\[
\varphi(t)=(1+t)^{p}-1-t^{p},\quad t\geqslant0,\;p>1.
\]
$\varphi(0)=0$，且 $\varphi'(t)=p\bigl[(1+t)^{\,p-1}-t^{\,p-1}\bigr]>0\;(t>0)$，故 $\varphi(t)>0\;(t>0)$。从而 $(A+B)^{p}>A^{p}+B^{p}$。两边开 $y$ 次方得
\[
(a^{x}+b^{x})^{\frac{1}{x}}=(A+B)^{\frac{1}{x}}>(A^{p}+B^{p})^{\frac{1}{xp}}=(a^{y}+b^{y})^{\frac{1}{y}}.
\]
	\end{enumerate}
	
	
}


\item 
(本题共 3 题，每小题 12 分，满分 36 分)
\begin{enumerate}
	%\renewcommand{\labelenumi}{\arabic{enumi}.}
	% A(\Alph) a(\alph) I(\Roman) i(\roman) 1(\arabic)
	%设定全局标号series=example	%引用全局变量resume=example
	%[topsep=-0.3em,parsep=-0.3em,itemsep=-0.3em,partopsep=-0.3em]
	%可使用leftmargin调整列表环境左边的空白长度 [leftmargin=0em]
	\item
求曲面 $x^{2}+y^{2}=a z$ 和 $z=2 a-\sqrt{x^{2}+y^{2}}(a>0)$ 所围成的立体体积。
	
	
	\item
	求级数 $\sum\limits_{n=1}^{\infty}\left[n(n+1)-\frac{1}{n(n+1)}\right] x^{n}$ 的和函数, 并求收敛域。
	
	
	\item
	求 $k$ 的取值范围, 使得方程 $\frac{k}{x}+x^{2}=1$ 有唯一正根。
	
	
	
	
\end{enumerate}



\banswer{
	\begin{enumerate}
		%\renewcommand{\labelenumi}{\arabic{enumi}.}
		% A(\Alph) a(\alph) I(\Roman) i(\roman) 1(\arabic)
		%设定全局标号series=example	%引用全局变量resume=example
		%[topsep=-0.3em,parsep=-0.3em,itemsep=-0.3em,partopsep=-0.3em]
		%可使用leftmargin调整列表环境左边的空白长度 [leftmargin=0em]
		\item
		$ V= \frac{ 5 }{ 6 } \pi a^{3} $
		\item 
		$\sum\limits_{n=1}^{\infty}\left[n(n+1)-\frac{1}{n(n+1)}\right] x^{n}=
		\left\{
		\begin{aligned}
			&0&(x=0)\\
			&\frac{2x}{(1-x)^{3}}-(1-x)\frac{\ln (1-x)}{x} -1 & (-1<x<1)
		\end{aligned}\right.
		$	
		\item 
		$ (-\infty, 0] \cup \{ \frac{2\sqrt{3}}{9} \} $
	\end{enumerate}
	
	
}


\newpage
\item 
(本题共 2 题，每题 12 分，满分 24 分)
\begin{enumerate}
	%\renewcommand{\labelenumi}{\arabic{enumi}.}
	% A(\Alph) a(\alph) I(\Roman) i(\roman) 1(\arabic)
	%设定全局标号series=example	%引用全局变量resume=example
	%[topsep=-0.3em,parsep=-0.3em,itemsep=-0.3em,partopsep=-0.3em]
	%可使用leftmargin调整列表环境左边的空白长度 [leftmargin=0em]
	\item
	将函数 $f(x)=\left\{\begin{array}{lc}x & -\pi \leqslant x \leqslant 0 \\ \pi x & 0<x \leqslant \pi\end{array}\right.$ 展开成傅里叶级数 (说明收敛情况), 并求$\sum\limits_{n=1}^{\infty} \frac{1}{n^{2}}$。
	
	
	\item 
确定常数 $\lambda$, 使得 $\frac{x}{y}\left(x^{2}+y^{2}\right)^{\lambda} d x-\frac{x^{2}}{y^{2}}\left(x^{2}+y^{2}\right)^{\lambda} d y=0$ 在 $D=\{(x, y) \mid y>0\}$ 内 为一全微分方程, 并利用曲线积分求此全微分方程的通解。
	
	
\end{enumerate}

\banswer{
\begin{enumerate}
	%\renewcommand{\labelenumi}{\arabic{enumi}.}
	\item
	$f(x)$ 可展开为
	\begin{align*}
	f(x)=\frac{\pi}{4}(\pi-1)+\sum_{n=1}^{\infty}\Bigl\{
	&-\frac{\pi-1}{n^{2}\pi}\left[1-(-1)^{n}\right]\cos nx\\
	&+\frac{\pi+1}{n}(-1)^{n+1}\sin nx\Bigr\} ,
	\end{align*}
	其和函数为
	\[
	\left\{\begin{array}{ll}
	x, & -\pi<x\leqslant0\\
	\pi x, & 0\leqslant x<\pi\\
	\frac{\pi(\pi-1)}{2}, & x=\pm\pi
	\end{array}\right.
	\]
	$f$ 在 $x=0$ 处左右极限都是 $0$，故级数在该点收敛于 $0$，于是
	\[
	\frac{\pi(\pi-1)}{4}=\frac{\pi-1}{\pi}\sum_{n=1}^{\infty}\frac{1-(-1)^{n}}{n^{2}}
	=\frac{2(\pi-1)}{\pi}\sum_{k=1}^{\infty}\frac{1}{(2k-1)^{2}} ,
	\]
	即 $\sum\limits_{k=1}^{\infty}\frac{1}{(2k-1)^{2}}=\frac{\pi^{2}}{8}$。再由
	$\sum\limits_{n=1}^{\infty}\frac{1}{n^{2}}
	=\sum\limits_{k=1}^{\infty}\frac{1}{(2k-1)^{2}}+\frac14\sum\limits_{n=1}^{\infty}\frac{1}{n^{2}}$
	解得
	\[
	\sum_{n=1}^{\infty}\frac{1}{n^{2}}=\frac43\cdot\frac{\pi^{2}}{8}=\frac{\pi^{2}}{6} .
	\]

	\par\vspace{1mm}\noindent
	\textbf{注：}{\kaishu
	原答案作 $\sum\limits_{n=1}^{\infty}\frac{1}{n^{2}}=\frac{\pi^{2}}{8}$，有误。
	由 $x=0$ 直接得到的是奇数项之和 $\frac{\pi^{2}}{8}\approx1.2337$，
	全体之和应为 $\frac{\pi^{2}}{6}\approx1.6449$（数值求和验证相符）。
	另原式中该因子处多写了一层左圆括号、致使左右定界符不配对，
	已改正为 $\left[1-(-1)^{n}\right]$。
	}\par

	\item
	$\lambda=-\frac12$，通解为 $\frac{\sqrt{x^{2}+y^{2}}}{y}=C$。
	记 $P=\frac{x}{y}\left(x^{2}+y^{2}\right)^{\lambda}$，
	$Q=-\frac{x^{2}}{y^{2}}\left(x^{2}+y^{2}\right)^{\lambda}$，则
	\begin{align*}
	\frac{\partial P}{\partial y}&=-\frac{x}{y^{2}}\left(x^{2}+y^{2}\right)^{\lambda}
	+2\lambda x\left(x^{2}+y^{2}\right)^{\lambda-1},\\
	\frac{\partial Q}{\partial x}&=-\frac{2x}{y^{2}}\left(x^{2}+y^{2}\right)^{\lambda}
	-\frac{2\lambda x^{3}}{y^{2}}\left(x^{2}+y^{2}\right)^{\lambda-1} .
	\end{align*}
	令二者相等，两端除以 $x\left(x^{2}+y^{2}\right)^{\lambda-1}$ 再乘 $y^{2}$，得
	\[
	-\left(x^{2}+y^{2}\right)+2\lambda y^{2}=-2\left(x^{2}+y^{2}\right)-2\lambda x^{2} ,
	\]
	即 $(1+2\lambda)\left(x^{2}+y^{2}\right)=0$，故 $\lambda=-\frac12$。此时
	\[
	P\,dx+Q\,dy=d\!\left(\frac{\sqrt{x^{2}+y^{2}}}{y}\right) ,
	\]
	（可沿 $D$ 内折线作曲线积分求得原函数），所以通解为 $\frac{\sqrt{x^{2}+y^{2}}}{y}=C$。

	\par\vspace{1mm}\noindent
	\textbf{注：}{\kaishu
	原答案作 $\lambda=\frac12$，有误。取 $\lambda=\frac12$ 时
	$\frac{\partial P}{\partial y}\neq\frac{\partial Q}{\partial x}$（数值校验：在 $(1.3,2.1)$ 处二者分别为
	$-0.2017$ 与 $-1.6578$），并非全微分方程；取 $\lambda=-\frac12$ 时二者相等。
	且原答案给出的通解 $\frac{\sqrt{x^{2}+y^{2}}}{y}=C$ 正对应 $\lambda=-\frac12$，
	与其自身所写的 $\lambda$ 并不自洽。
	}\par
\end{enumerate}
}


	
\end{enumerate}


